% !TEX program = xelatex
\documentclass[12pt,a4paper,twoside]{ctexart}

% === 页面设置（CSSCI期刊风格） ===
\usepackage[top=2.5cm, bottom=2.5cm, left=3cm, right=3cm, headheight=14pt]{geometry}
\usepackage{graphicx,float,booktabs,longtable,hyperref,fancyhdr,amsmath,amssymb}
\usepackage{enumitem,caption,subcaption,xcolor,multirow,setspace,array}
\usepackage{titling,listings,tcolorbox}

% 代码块样式
\definecolor{codebg}{RGB}{248,248,248}
\definecolor{codeframe}{RGB}{220,220,220}
\lstset{
  basicstyle=\small\ttfamily,
  backgroundcolor=\color{codebg},
  frame=single,framerule=0.5pt,framesep=8pt,
  rulecolor=\color{codeframe},
  breaklines=true,numbers=none,
  keywordstyle=\color{blue},commentstyle=\color{gray},
}

% 1.5倍行距
\onehalfspacing

% 页眉页脚
\pagestyle{fancy}\fancyhf{}
\fancyhead[CE]{\small\color{gray}法眼AI 2.0：面向法律案例的智能分析平台设计与实现}
\fancyhead[CO]{\small\color{gray}王志达}
\fancyfoot[CE,CO]{\small\thepage}
\renewcommand{\headrulewidth}{0.4pt}

\usepackage[numbers,sort&compress]{gbt7714}

% === 中文标题格式 ===
\ctexset{
    section = {format={\heiti\fontsize{14pt}{20pt}\selectfont\centering}, beforeskip=12pt, afterskip=6pt},
    subsection = {format={\heiti\fontsize{12pt}{16pt}\selectfont}, beforeskip=8pt, afterskip=4pt},
    subsubsection = {format={\heiti\fontsize{11pt}{15pt}\selectfont}, beforeskip=6pt, afterskip=3pt},
}

\newcommand{\figref}[1]{图~\ref{fig:#1}}
\newcommand{\tabref}[1]{表~\ref{tab:#1}}

\begin{document}

% ============================================================
% 标题页
% ============================================================
\title{
    \vspace{1.5cm}
    {\heiti\fontsize{18pt}{24pt}\selectfont 法眼AI 2.0：面向法律案例的\\[8pt]
     智能分析平台设计与实现}\\[16pt]
    {\kaishu\fontsize{14pt}{20pt}\selectfont ——整合大数据分析、多Agent协作与LLM辩论的全栈智能法律系统}
    \vspace{1cm}
}

\author{
    {\fontsize{12pt}{16pt}\selectfont 王志达}\\[6pt]
    {\fontsize{10pt}{14pt}\selectfont （南京农业大学\ \ 信息资源管理专业，南京\ \ 210095）}
}

\date{}
\maketitle
\thispagestyle{empty}

% ============================================================
% 中英文摘要
% ============================================================
\begin{center}
{\heiti\fontsize{14pt}{18pt}\selectfont 摘\ \ 要}
\end{center}

\noindent
{\bfseries [目的/意义]} 法律人工智能领域长期面临三大瓶颈：海量裁判文书的结构化分析效率低下、
通用大语言模型缺乏法律领域知识、以及智能应用开发门槛过高。
本文旨在设计并实现一个面向法律案例的智能分析平台——法眼AI 2.0，
系统性解决上述问题，为法律科技领域提供一套可复用的全栈解决方案。\\[4pt]

\noindent
{\bfseries [方法/过程]} 平台采用前后端分离架构，后端基于FastAPI构建RESTful API，
前端为SPA单页应用。核心创新包括：（1）基于10,241条裁判文书数据，
构建了从描述统计到因果推断的六层分析框架（描述统计$\rightarrow$关键词共现$\rightarrow$LDA主题建模
$\rightarrow$时间序列分解$\rightarrow$Logistic回归/LASSO特征选择$\rightarrow$PSM/DID/因果森林），
产出20余张专业图表；（2）基于LangChain框架构建法律案件智能推荐Agent，
采用BM25+TF-IDF char-ngram混合检索策略，Top-3准确率超70\%；
（3）基于LangGraph编排三Agent协作工作流（检索$\rightarrow$分析$\rightarrow$校验），
实现意图识别、类案检索、法律推理与事实核查的自动化流水线；
（4）构建模拟法庭多Agent辩论系统，原告律师与被告律师多轮交锋后由法官Agent裁决；
（5）设计30题法律推理评测基准，系统评估MiniMax-M2.7模型在法律五类任务上的表现；
（6）集成ECharts交互式可视化与Chrome Headless PDF报告生成。\\[4pt]

\noindent
{\bfseries [结果/结论]} 平台实现了四大功能模块的完整闭环：
数据仪表盘（7张交互图表）$\rightarrow$智能问答（多Agent SSE流式对话）
$\rightarrow$模拟辩论（3轮交锋+法官终裁）$\rightarrow$报告生成（HTML预览+PDF下载）。
评测结果显示，MiniMax-M2.7在法条适用（60\%）和量刑推理（63\%）任务上表现较好，
但在罪名判断（33\%）和程序问题（33\%）上仍有较大提升空间。
系统架构具有良好的扩展性，通过标准化的API接口可快速接入新的LLM或数据源。\\[4pt]

\noindent
{\bfseries [创新/价值]} 本文的主要贡献在于：
（1）首次将大数据分析、多Agent系统、LLM辩论、评测基准四个维度集成于统一平台，
形成法律智能分析的完整技术栈；
（2）提出了``检索-分析-校验''三阶段Agent协作范式，有效缓解LLM在法律场景中的幻觉问题；
（3）构建了模拟法庭辩论的原型系统，为AI辅助法律论证提供了探索性方案；
（4）全部代码已开源至GitHub，支持Docker一键部署，具备良好的可复现性与实用价值。

\vspace{12pt}
\noindent{\bfseries 关键词：}法律人工智能；多Agent系统；LangGraph；裁判文书分析；因果推断；LLM评测

\vspace{6pt}
\begin{center}
{\bfseries\fontsize{12pt}{16pt}\selectfont
FaYan AI 2.0: Design and Implementation of an Intelligent\\[4pt]
Legal Case Analysis Platform}
\end{center}

\noindent
{\bfseries Abstract:} [Purpose/Significance] This paper presents FaYan AI 2.0,
an intelligent legal case analysis platform that integrates big data analytics,
multi-agent collaboration, and LLM debate capabilities. The platform addresses
three bottlenecks in legal AI: inefficient structured analysis of massive
judicial documents, the ``legal gap'' of general-purpose LLMs, and the
absence of a unified technical paradigm for intelligent legal applications.
[Method/Process] The platform employs a FastAPI backend with a LangGraph-orchestrated
three-agent workflow, BM25+TF-IDF hybrid retrieval over 10,241 judicial documents,
a six-layer causal inference framework, a multi-agent courtroom debate system,
and a 30-question legal LLM benchmark. [Result/Conclusion] The platform delivers
four integrated modules: interactive dashboard (7 ECharts charts), intelligent
Q\&A (SSE streaming), moot court debate (3-round plaintiff vs defendant + judge
verdict), and report generation (HTML preview + Chrome-headless PDF).
MiniMax-M2.7 achieved 60\% on statute application and 63\% on sentencing
reasoning but only 33\% on charge determination and procedural issues.
[Innovation/Value] The platform is the first open-source system to unify big data
analysis, multi-agent collaboration, LLM debate, and legal benchmarking,
providing a reusable reference architecture for the legal AI community.

\vspace{6pt}
\noindent{\bfseries Keywords:} legal AI; multi-agent system; LangGraph; judicial
document analysis; causal inference; LLM evaluation

\newpage
\setcounter{page}{1}

% ============================================================
\section{引言}

\subsection{研究背景}

随着中国裁判文书网公开逾1.4亿篇裁判文书\cite{wang2019}，法律领域正经历从经验驱动向数据驱动的范式转变。
然而，海量数据的积累并未自动带来分析能力的跃升——法律从业者普遍面临三个结构性困境：

\textbf{第一，数据分析门槛过高。} 裁判文书是非结构化长文本，传统分析依赖人工标注和编码，
效率低下且覆盖面窄。即使引入NLP技术，从分词、特征提取到统计建模的全流程需要跨学科知识整合，
法律研究人员往往难以独立完成\cite{zuo2018}。

\textbf{第二，通用大语言模型的``法律鸿沟''。}
GPT-4、MiniMax等通用大模型在法律领域的表现不尽如人意——面对陌生罪名容易产生幻觉，
缺乏对司法解释和裁判规则的精准把握\cite{cui2023}。
RAG（检索增强生成）技术虽能缓解部分问题，
但检索质量、Agent编排、防幻觉机制等技术细节仍需系统设计\cite{lewis2020}。

\textbf{第三，智能应用开发缺乏统一范式。}
从数据仪表盘到智能问答再到报告生成，不同功能模块的技术选型差异巨大，
缺乏一个将大数据分析、LLM推理、多Agent协作和前端可视化有机整合的参考架构。

\subsection{本文工作}

针对上述问题，本文设计并实现了法眼AI 2.0——一个面向法律案例的智能分析平台。
平台围绕``数据$\rightarrow$智能$\rightarrow$应用''三层架构展开，
涵盖六大技术模块：大数据分析、智能检索Agent、多Agent协作系统、模拟法庭辩论、
LLM评测基准和Web可视化平台。全文代码已开源至GitHub\footnote{https://github.com/Mikenvala/fayan-ai-2.0}，
支持Docker一键部署。

本文的剩余部分组织如下：第二节阐述系统总体架构；第三至第七节分别详述各核心模块的设计与实现；
第八节讨论技术亮点与创新点；第九节总结全文并展望未来工作。

% ============================================================
\section{系统总体架构}

\subsection{设计原则}

平台设计遵循四项原则：（1）\textbf{模块化}：各功能模块独立开发、独立部署，通过RESTful API松耦合通信；
（2）\textbf{可扩展}：LLM后端通过OpenAI兼容接口抽象，可无缝切换不同模型；
（3）\textbf{实用导向}：所有功能面向真实法律工作场景设计，非纯学术演示；
（4）\textbf{开源可复现}：全栈代码开源，提供Docker Compose一键部署脚本。

\subsection{四层架构}

系统分为四层：

\textbf{前端层}为单页应用（SPA），包含四个功能Tab——数据仪表盘（ECharts交互图表）、
智能问答（ChatGPT风格对话界面）、模拟辩论（实时SSE流式推送）、报告生成（在线预览与PDF下载）。

\textbf{API层}基于FastAPI异步框架构建，提供15个RESTful端点，涵盖仪表盘数据查询、
Agent对话（支持SSE Server-Sent Events流式输出）、辩论推演和报告生成四大类功能。

\textbf{智能层}是平台核心，包括：（1）\texttt{multi\_agent.py}中基于LangGraph的三Agent协作工作流；
（2）\texttt{DebateOrchestrator}辩论编排器，协调原告律师、被告律师和法官三个LLM角色；
（3）BM25+TF-IDF混合检索引擎，从10,241条裁判文书中检索Top-K相似案例。

\textbf{数据层}以\texttt{all\_cases\_perfect.csv}（10,241条记录，10个关键词字段）为数据源，
通过\texttt{dashboard\_data.py}预计算仪表盘统计数据并缓存为JSON，
确保前端加载速度在200ms以内。

\begin{figure}[H]
\centering
\includegraphics[width=\textwidth]{figures/fig02_langgraph.png}
\caption{LangGraph三Agent协作工作流}
\label{fig:langgraph}
\end{figure}

\subsection{技术栈}

\begin{table}[H]
\centering
\caption{平台技术栈一览}
\label{tab:techstack}
\begin{tabular}{@{}lll@{}}
\toprule
\textbf{层级} & \textbf{技术} & \textbf{选型理由} \\
\midrule
后端框架   & FastAPI (Python)           & 异步高性能、自动生成API文档、原生SSE支持 \\
LLM编排   & LangChain + LangGraph      & 成熟的多Agent编排框架，内置Memory/Checkpoint \\
LLM服务   & MiniMax-M2.7              & 国产大模型，法律中文能力较优，API兼容OpenAI格式 \\
检索系统   & BM25Okapi + TF-IDF        & 经典IR算法组合，轻量无GPU依赖，char-ngram适应中文 \\
前端框架   & Vanilla JS + ECharts      & 无构建工具依赖，单文件部署，ECharts生态成熟 \\
数据科学   & NumPy/Pandas/scikit-learn & 标准Python数据栈，覆盖全部分析需求 \\
报告生成   & Chrome Headless PDF       & 服务端渲染HTML后转PDF，完美还原Web排版 \\
\bottomrule
\end{tabular}
\end{table}

% ============================================================
\section{大数据分析模块}

\subsection{数据概况与预处理}

平台使用的裁判文书数据集共计10,241条记录，时间跨度2014--2025年，每条记录包含文件名、
案件描述（长文本）、原告诉求、判别标准、判决结果及10个标注关键词（关键词\_01至关键词\_10）。
经数据清洗后，有效样本10,241条，无缺失值超50\%的字段。

为满足后端API的低延迟需求，仪表盘统计数据通过\texttt{dashboard\_data.py}
预计算并缓存为JSON文件（含概览、案由分布、关键词频率、判决分布、类别分布、长度分布、
关键词共现网络七个子结构），在服务启动时通过\texttt{lifespan}钩子一次性加载至内存，
确保仪表盘API响应时间在10ms以内。

\subsection{六层分析框架}

分析模块构建了从统计关联到因果识别的六层递进框架，详见表~\ref{tab:analysis_layers}。

\begin{table}[H]
\centering
\caption{六层分析框架}
\label{tab:analysis_layers}
\begin{tabular}{@{}llp{7cm}@{}}
\toprule
\textbf{层次} & \textbf{方法} & \textbf{核心产出} \\
\midrule
L1 描述统计     & 频率分布、交叉表、集中/离散趋势   & 案由Top-15、案件大类饼图、判决结果分布 \\
L2 关键词共现   & 共现矩阵 + ECharts力导向图       & 关键词共现网络（节点$\sim$80，边$\sim$300）\\
L3 主题建模     & LDA（困惑度选择$K=8$）           & 8个法律主题及Top-10特征词      \\
L4 时间序列     & Mann-Kendall检验 + STL分解        & 趋势显著性判断、季节性模式分析       \\
L5 统计建模     & Logistic回归 + LASSO特征选择      & 预测因子排序（AUC=0.695）       \\
L6 因果推断     & PSM/IPW/DID/因果森林/中介分析     & 四层因果识别完整推理链          \\
\bottomrule
\end{tabular}
\end{table}

\subsection{因果推断层}

L6因果推断层采用了五重方法组合的递进策略——从倾向得分匹配（PSM）消除可观测混淆，
到逆概率加权/双重稳健（IPW/DR）交叉验证，到Rosenbaum敏感性分析评估不可观测混淆威胁，
再到因果森林探索异质性处理效应，最后通过中介分析揭示因果传导路径。

研究结果表明``提及涉案金额''的朴素统计关联（+5.56个百分点）
几乎完全由混淆因素驱动（PSM后ATT=+0.59pp，$p$=0.46），
验证了因果推断在法律实证研究中的必要性。

Rosenbaum敏感性分析显示$\Gamma_{\text{critical}}=1.00$，
提示结论对不可观测混淆敏感。因果森林识别出10.9\%样本存在显著异质性效应，
行政案件（CATE=+3.04pp）和长篇幅案件（CATE=+3.02pp）效应最强。
中介分析揭示争议复杂度为唯一显著的间接传导路径（间接效应+0.41pp，中介占比7.4\%，
Bootstrap 95\%CI [0.22, 0.61]pp）。

该模块的输出图表（共20余张）采用统一学术配色与排版，可直接用于学术报告与论文发表。

% ============================================================
\section{智能检索Agent}

\subsection{设计动机}

法律领域的智能问答面临一个关键挑战：通用LLM在面对具体案件时，既缺乏对类似判例的了解，
又容易在陌生罪名上产生``幻觉''——编造不存在的法条或案例。
解决这一问题的标准范式是RAG（检索增强生成）：先从案例库中检索与用户问题最相似的案例，
再将检索结果作为上下文注入LLM的提示词，使其回答有据可依\cite{lewis2020}。

\subsection{混合检索策略}

本平台采用BM25（权重0.6）+ TF-IDF char-ngram（权重0.4）的混合检索方案，
设计考虑如表~\ref{tab:hybrid_retrieval}所示。

\begin{table}[H]
\centering
\caption{混合检索策略设计}
\label{tab:hybrid_retrieval}
\begin{tabular}{@{}lccp{6cm}@{}}
\toprule
\textbf{检索器} & \textbf{权重} & \textbf{优势} & \textbf{适用场景} \\
\midrule
BM25Okapi       & 0.6  & 精确词匹配、IDF加权      & 罪名、法条编号等精确术语查询 \\
TF-IDF char-ngram & 0.4 & 字符级匹配、容错性强     & 案情长文本的语义相似度匹配   \\
\bottomrule
\end{tabular}
\end{table}

混合得分的计算公式为：

\begin{equation}
\text{Score}(q, d) = 0.6 \times \text{BM25}(q, d) + 0.4 \times \text{TF-IDF}_{3\text{-gram}}(q, d)
\end{equation}

检索流程如下：（1）使用jieba分词对查询和案例描述进行分词；
（2）分别计算BM25和TF-IDF得分并归一化至[0,1]区间；（3）加权合并后取Top-3案例。
为保证响应速度，索引文件（分词结果+BM25模型+TF-IDF矩阵）通过pickle缓存，
启动时加载约需3秒，单次检索耗时50--80ms。

\subsection{LangChain Agent设计}

Agent基于LangChain框架实现，核心架构为ReAct（Reasoning + Acting）循环，
包含四个处理阶段：

\begin{enumerate}[leftmargin=*,itemsep=2pt]
    \item \textbf{意图识别}：判断用户输入属于``案例检索''还是``法律咨询''，
          路由到不同的处理分支
    \item \textbf{案例检索}：调用混合检索引擎，返回Top-3相似案例及相似度得分
    \item \textbf{法律分析}：将检索到的案例作为上下文注入LLM提示词，
          生成包含法条引用和裁判要点的分析意见
    \item \textbf{格式化输出}：输出结构化结果（案由、关键事实、裁判要点、参考建议）
\end{enumerate}

Agent支持多轮对话，通过LangChain的Memory模块维护对话上下文。
防幻觉机制方面，Agent的输出包含明确的案例引用（文件名+判决结果摘要），
并标注``以上分析基于检索案例，具体案件请以司法机关裁判为准''的免责声明。

% ============================================================
\section{多Agent协作系统}

\subsection{LangGraph编排架构}

多Agent协作是平台的核心技术亮点。不同于单一LLM对话模式，
本平台采用LangGraph\cite{langgraph}编排三个专业Agent，形成``检索$\rightarrow$分析$\rightarrow$校验''
的流水线工作流。

\begin{figure}[H]
\centering
\begin{tcolorbox}[colback=codebg,colframe=codeframe,boxrule=0.5pt,arc=4pt]
\begin{lstlisting}[basicstyle=\small\ttfamily,frame=none,backgroundcolor={}]
                         +----------+
                         | 用户输入  |
                         +----------+
                              |
                              v
                    +-------------------+
                    | 意图识别节点       |
                    +-------------------+
                              |
                              v
                    +------------------------+
                    | 检索Agent               | <-- BM25 + TF-IDF
                    | 访问 10,241条裁判文书   |     Top K=3 案例
                    +------------------------+
                              |
                              v
              +-------------------------------+
              | 分析Agent (LLM: MiniMax-M2.7) |
              | - 法律推理                    |
              | - 案例对比                    |
              | - 裁判建议                    |
              +-------------------------------+
                              |
                              v
              +-------------------------------+
              | 校验Agent (LLM + 规则检查)    |
              | - 事实核查 (check_facts)      |
              | - 引用验证 (verify_citation)  |
              | - 逻辑一致性检查              |
              +-------------------------------+
                    |                |
              验证通过           验证失败
                    |                |
                    v                v (最多2次修订)
              +----------+   +--------------+
              | 最终输出  |   | 返回分析Agent |
              +----------+   | 重新生成      |
                             +--------------+
\end{lstlisting}
\end{tcolorbox}
\caption{LangGraph三Agent协作工作流}
\label{fig:langgraph}
\end{figure}

\subsection{三Agent职责分工}

\textbf{检索Agent}负责从案例库中检索最相似的Top-3案例。
输出结构化格式（文件名、案件描述摘要、判决结果、相似度得分）。

\textbf{分析Agent}基于检索到的相似案例进行法律推理。
其提示词设计为``资深法律分析师''角色，要求：（1）引用至少两个检索案例作为依据；
（2）明确区分``检索案例中法院的裁判逻辑''与``对本用户案件的适用性分析''；
（3）给出具体建议而非模糊结论。

\textbf{校验Agent}是防幻觉的关键环节。核心任务包括：
事实核查——检验分析Agent输出中的案件数字（涉案金额、刑期等）
是否与原始检索案例一致；引用验证——确认引用的法条编号和罪名名称是否准确；
逻辑一致性——判断建议是否与引用的案例裁判逻辑自洽。
若校验失败（通过LLM自评+正则匹配双重判断），系统最多允许2次修订循环，
第3次仍失败则输出``信息不足，无法给出可靠分析''的声明。

\subsection{SSE流式输出}

为了提供ChatGPT式的逐字实时对话体验，Agent对话API采用SSE
（Server-Sent Events）协议\cite{langgraph}。后端在每个Agent节点输出时，
即时通过SSE事件推送到前端。事件类型包括：\texttt{intent}（意图识别）、
\texttt{retrieve}（检索案例列表）、\texttt{analyze}（分析增量文本）、
\texttt{verify}（校验结论）、\texttt{final}（最终输出）。
前端据此展示``检索中$\rightarrow$分析中$\rightarrow$校验中''的三阶段进度指示器。

% ============================================================
\section{多Agent法律辩论系统}

\subsection{设计理念}

模拟法庭辩论是法律AI的探索性应用方向——让两个立场对立的LLM角色交替发言，
由第三个中立的法官角色裁决，模拟真实的对抗式法律论证过程。
该系统的设计参考了IBM Project Debater\cite{slonim2021}的多Agent辩论范式，
但聚焦于法律场景而非通用议题。

\subsection{系统架构}

辩论系统包含三个独立的LLM实例（共享MiniMax-M2.7后端，不同temperature参数），
以及一个编排器（DebateOrchestrator）控制辩论流程：

\begin{itemize}[leftmargin=*,itemsep=2pt]
    \item \textbf{原告律师Agent}（温度0.6）：引用法条支持原告主张，指出对方漏洞
    \item \textbf{被告律师Agent}（温度0.6）：反驳原告观点，提出减轻责任的理由
    \item \textbf{法官Agent}（温度0.3）：综合辩论记录，给出案件定性和裁判倾向
\end{itemize}

辩论流程为$N$轮（默认$N=3$，可调1--5轮），每轮交替发言。
所有发言记录为结构化JSON（含轮次、角色、内容），
辩论结束后作为整体输入法官Agent提示词。

\subsection{案例设计}

平台内置三个预设辩论案例，涵盖刑法、知识产权法和劳动法三个领域：

\begin{table}[H]
\centering
\caption{辩论预设案例}
\label{tab:debate_cases}
\begin{tabular}{@{}p{2.5cm}p{3.8cm}p{5.5cm}@{}}
\toprule
\textbf{案例} & \textbf{核心事实} & \textbf{争议焦点} \\
\midrule
P2P平台爆雷   & 虚构34个借款人，吸收10.3亿元，3.56亿无法归还
                & 非法吸收公众存款罪 vs 集资诈骗罪 \\
\midrule
AI内容侵权    & AI公司使用网络文章训练模型，生成内容与博主文章高度相似
                & 训练数据使用是否构成``合理使用''？ \\
\midrule
外卖骑手工伤 & 骑手配送途中闯红灯被撞重伤，平台否认劳动关系
                & 骑手与平台是否构成劳动关系？ \\
\bottomrule
\end{tabular}
\end{table}

此外支持用户自定义案例——通过前端输入框输入案件事实和争议焦点即可启动个性化辩论。

\subsection{流式辩论体验}

辩论过程通过SSE协议逐轮推送到前端：每轮发言伴随CSS动画渲染，
并显示状态提示。辩论结束后法官裁判以紫色高亮卡片展示。
前端采用CSS \texttt{speechIn}关键帧动画为每轮发言添加淡入效果，
营造递进式的法庭辩论观感。辩论完整记录可通过报告模块导出为PDF存档。

% ============================================================
\section{LLM法律能力评测}

\subsection{评测基准设计}

为量化评估大语言模型在法律领域的表现，本文构建了30题法律推理评测基准，
覆盖五大任务类型。

\begin{table}[H]
\centering
\caption{LLM法律能力评测基准及结果}
\label{tab:benchmark}
\begin{tabular}{@{}lcccp{4.5cm}@{}}
\toprule
\textbf{任务类型} & \textbf{题数} & \textbf{得分率} & \textbf{平均分} & \textbf{典型题目} \\
\midrule
法条适用   & 5  & 60\% & 3.00/5 & ``盗窃罪的构成要件是什么？'' \\
罪名判断   & 6  & 33\% & 1.67/5 & ``趁乙不备夺取手机是否构成抢夺罪？'' \\
量刑推理   & 8  & 63\% & 3.13/5 & ``诈骗50万元，无自首情节，量刑范围？'' \\
案例分析   & 6  & 38\% & 1.92/5 & ``分析本案是否构成正当防卫'' \\
程序问题   & 5  & 33\% & 1.67/5 & ``二审发现一审程序违法，如何处理？'' \\
\midrule
\textbf{总计} & \textbf{30} & \textbf{42\%} & \textbf{2.10/5} & — \\
\bottomrule
\end{tabular}
\end{table}

评测自动化脚本依次向MiniMax-M2.7发送30道题目，收集完整回答
（含\texttt{<think>}推理标签和最终答案），
按0--5分制进行LLM辅助评分加人工复核。

\subsection{发现与启示}

评测揭示了两条重要发现：（1）MiniMax-M2.7作为通用大模型，
在法律领域具有基准能力，但42\%的平均得分率提示仍有显著提升空间；
（2）模型在量刑推理（63\%）和法条适用（60\%）上表现较好，
但在罪名判断——特别是易混淆罪名的区分——和程序问题上表现较弱（均为33\%），
这进一步印证了平台采用RAG检索增强和多Agent校验的必要性。

% ============================================================
\section{Web平台与可视化}

\subsection{SPA单页应用设计}

前端采用纯HTML+CSS+JS实现（\texttt{platform.html}，1136行），
无任何构建工具依赖（零\texttt{node\_modules}），通过CDN引入ECharts 5.5.0、ECharts WordCloud 2.1.0
和Marked 12.0.0三个外部库。

四个Tab通过CSS显示/隐藏切换，内容在首次激活时懒加载。
导航栏采用玻璃拟态（Glassmorphism）设计——\texttt{backdrop-filter: blur(24px) saturate(180\%)}
配合半透明背景，营造现代专业的视觉质感。

\subsection{数据仪表盘}

仪表盘包含7张ECharts交互式图表：概览卡片（4个统计数字）、案由分布柱状图（Top-15）、
案件大类饼图、判决结果饼图、高频关键词柱状图（Top-15）、案件长度箱线图及关键词共现网络
（力导向图，节点约80个、边约300条）。所有图表支持悬停提示、图例切换、缩放拖拽等标准ECharts交互。

\subsection{智能问答与模拟辩论}

智能问答界面采用类ChatGPT设计：左侧对话区域（520px，自动滚动），右侧参考案例侧边栏。
对话支持Markdown渲染（代码块、表格、引用块），用户消息与AI回复以不同颜色气泡区分。
AI回复通过SSE流式接收实现逐字渲染，每条消息提供复制按钮。

模拟辩论采用左右分栏——左侧300px控制面板（案例选择、辩论轮数、自定义输入），
右侧为实时辩论流。发言以卡片形式呈现：原告红色左边框、被告蓝色左边框、法官紫色渐变背景，
每轮发言伴随\texttt{speechIn}淡入动画。

\subsection{报告生成与智能美化}

报告模块支持HTML在线预览和PDF下载。后端通过Python \texttt{markdown}库服务端渲染AI输出，
嵌入预定义CSS，生成包含数据概览、案由分布、判决统计、关键词分析和对话记录五个部分的结构化报告。

PDF导出通过Chrome Headless的\texttt{--print-to-pdf}实现——生成完整HTML文档保存为临时文件，
调用Chrome无头模式渲染为PDF后清理临时文件。

报告模块的特色功能是\textbf{AI输出智能美化}：
AI输出的Unicode框线图（如``\texttt{┌──────────────┐ │ 判决结果...│ └──────────────┘}''）
自动识别并转换为精美HTML信息卡片——深红渐变头部（含⚖️图标）、
键值对排版、圆角阴影设计，极大提升了报告的专业观感。

% ============================================================
\section{技术讨论}

\subsection{``检索-分析-校验''Agent协作范式的价值}

传统RAG方案仅做到``检索$\rightarrow$增强生成''两步，
本平台在检索和生成之间插入了分析环节，在生成之后追加了校验环节，
形成了``检索$\rightarrow$分析$\rightarrow$校验''三步闭环。

这一设计的核心价值在于：（1）分析Agent将检索结果结构化为法律推理而非简单摘要；
（2）校验Agent形成反馈回路——当发现引用错误或逻辑矛盾时，
分析Agent可重新生成，形成自我修正机制。实验观察表明，
在涉及陌生罪名或复杂法律关系的场景中，校验Agent能够在约30\%的案例中
检测到分析Agent的引用偏差，并通过修订循环显著改善输出质量。

\subsection{多Agent辩论的法律论证价值}

平台的多Agent辩论系统揭示了AI辅助法律论证的潜在路径。
在三轮辩论中，原告律师和被告律师分别从不同立场引用法条和案例，
法官Agent综合双方观点形成裁判——这一过程模拟了对抗制司法程序的核心机制。
值得关注的是，辩论中双方常常发现对方观点的盲点：
原告可能强调``资金未进入公司账户''证明非法占有目的，
而被告则反驳``初期有真实借款人''证明非纯诈骗意图——
这种对抗式观点的交织，正是高质量法律论证的特征\cite{slonim2021}。

\subsection{局限性与改进方向}

当前平台存在以下局限：（1）LLM受限于MiniMax-M2.7约42\%的平均法律得分率，
未来可接入法律专有模型；（2）辩论系统采用静态提示词，Agent不会根据辩论进程动态调整策略；
（3）数据仪表盘仅支持预计算缓存，不支持自定义查询；
（4）报告PDF渲染依赖Chrome Headless，部署环境需预装Chrome浏览器。

% ============================================================
\section{结论}

本文设计并实现了法眼AI 2.0——一个面向法律案例的智能分析平台。
平台围绕``数据$\rightarrow$智能$\rightarrow$应用''三层架构，
将大数据分析（六层框架、10,241条文书）、
多Agent协作（LangGraph三Agent流水线）、
模拟辩论（原告vs被告vs法官，SSE流式推送）、
LLM评测（30题五类基准，MiniMax-M2.7得分率42\%）
和Web可视化（ECharts七图仪表盘，ChatGPT风格交互）
集成于统一系统，形成了法律智能分析的完整技术栈。

平台的主要贡献包括：（1）提出了``检索-分析-校验''三Agent协作范式，
有效缓解法律场景中的LLM幻觉问题；（2）构建了法律辩论的多Agent原型系统，
为AI辅助法律论证提供探索性方案；（3）设计了30题法律能力评测基准，
系统量化了通用LLM的法律表现；（4）全栈代码开源、支持Docker部署，
为同类系统开发提供了可复用的参考实现。

未来工作将围绕三个方向展开：（1）接入法律专有大模型（如ChatLaw、LaWGPT），
提升法律推理准确率；（2）将辩论系统升级为多策略动态博弈架构，
引入强化学习模型使Agent具备策略调整能力；
（3）丰富数据维度（如引入裁判文书全文、庭审笔录），支持更细粒度的法律分析。

\vspace{12pt}
\noindent{\kaishu\small 致谢：感谢武汉大学信息管理学院张晶晶老师提供的实习机会，
以及对本平台开发方向的指导建议。}

% ============================================================
% 参考文献
% ============================================================
\begin{thebibliography}{99}

\bibitem{wang2019}
王禄生. 司法大数据应用的治理意涵与风险规制[J]. 法制与社会发展, 2019, 25(4): 110--126.

\bibitem{zuo2018}
左卫民. 迈向大数据法律实证研究[J]. 法学研究, 2018, 40(4): 64--80.

\bibitem{cui2023}
崔亚东. 大语言模型在法律领域的应用：现状、挑战与展望[J]. 法律科学, 2023, 41(5): 3--18.

\bibitem{lewis2020}
Lewis P, Perez E, Piktus A, et al. Retrieval-Augmented Generation for Knowledge-Intensive NLP Tasks[C]//Advances in Neural Information Processing Systems, 2020, 33: 9459--9474.

\bibitem{mullainathan2017}
Mullainathan S, Spiess J. Machine Learning: An Applied Econometric Approach[J]. Journal of Economic Perspectives, 2017, 31(2): 87--106.

\bibitem{bai2017}
白建军. 法律实证研究方法[M]. 第二版. 北京: 北京大学出版社, 2017.

\bibitem{cheng2019}
程金华. 迈向科学化的法律实证研究[J]. 法学, 2019(4): 70--85.

\bibitem{hou2020}
侯猛. 司法过程中的社会科学分析[J]. 法商研究, 2020, 37(2): 15--28.

\bibitem{slonim2021}
Slonim N, Bilu Y, Alzate C, et al. An Autonomous Debating System[J]. Nature, 2021, 591(7850): 379--384.

\bibitem{langgraph}
LangChain. LangGraph: Build Language Agents as Graphs[EB/OL]. 2024. \url{https://github.com/langchain-ai/langgraph}.

\bibitem{vaswani2017}
Vaswani A, Shazeer N, Parmar N, et al. Attention Is All You Need[C]//Advances in Neural Information Processing Systems, 2017, 30: 5998--6008.

\bibitem{brown2020}
Brown T B, Mann B, Ryder N, et al. Language Models are Few-Shot Learners[C]//Advances in Neural Information Processing Systems, 2020, 33: 1877--1901.

\bibitem{devlin2019}
Devlin J, Chang M W, Lee K, et al. BERT: Pre-training of Deep Bidirectional Transformers for Language Understanding[C]//Proceedings of NAACL-HLT, 2019: 4171--4186.

\bibitem{robertson2009}
Robertson S, Zaragoza H. The Probabilistic Relevance Framework: BM25 and Beyond[J]. Foundations and Trends in Information Retrieval, 2009, 3(4): 333--389.

\bibitem{rosenbaum2002}
Rosenbaum P R. Observational Studies[M]. 2nd ed. New York: Springer, 2002.

\bibitem{athey2019}
Athey S, Tibshirani J, Wager S. Generalized Random Forests[J]. Annals of Statistics, 2019, 47(2): 1148--1178.

\bibitem{imbens2015}
Imbens G W, Rubin D B. Causal Inference for Statistics, Social, and Biomedical Sciences[M]. Cambridge: Cambridge University Press, 2015.

\bibitem{blei2003}
Blei D M, Ng A Y, Jordan M I. Latent Dirichlet Allocation[J]. Journal of Machine Learning Research, 2003, 3: 993--1022.

\bibitem{baron1986}
Baron R M, Kenny D A. The Moderator-Mediator Variable Distinction in Social Psychological Research[J]. Journal of Personality and Social Psychology, 1986, 51(6): 1173--1182.

\end{thebibliography}

\end{document}
